\chapter{The 2-Basis Projection Framework}
\label{ch:bridging_metrology_imaging}

A central contribution of this thesis is the formal mathematical and physical unification of two seemingly distinct tracks in quantum information science: the variational estimation of mixed-state metrological bounds on Noisy Intermediate-Scale Quantum (NISQ) simulators (Chapter~\ref{ch:methods}) and spatial-mode demultiplexing (SPADE) for sub-Rayleigh imaging and spectroscopy of highly unbalanced optical sources (Chapter~\ref{sec:astrophysical_metrology_spade}).

While Chapter~\ref{ch:methods} established the abstract theory of the Variational Quantum State Eigensolver (VQSE) and the Variational Quantum Fisher Information Estimation (VQFIE) protocols to resolve classical diagonalization bottlenecks, and Chapter~\ref{sec:astrophysical_metrology_spade} motivated the physical and astrophysical challenge of sub-diffraction telescopy, this chapter establishes the rigorous bridge between these two worlds.

Rather than treating them as isolated tracks, we prove that the physical apparatus of spatial-mode demultiplexing is a physical, analog realization of the rank-truncation projector $\Pi^{(1)}$ and the Truncated Quantum Fisher Information (TQFI) formalism developed in our variational metrology pipeline. We demonstrate how telescope alignment corresponds directly to VQSE majorization, how the ultimate limits of SPADE saturate the variational TQFI lower bounds, how low-rank mode-sorting is justified by Krylov subspace closure, and how physics-informed neural network decoders act as the physical analogue of isometric polar completions. Finally, we analyze the global precision limits via the joint Quantum Fisher Information Matrix (QFIM) to prove that alignment errors and parameter estimates completely decouple as $\epsilon \to 0$.

% ==============================================================================
\section{SPADE as a Physical Rank-Truncation Projector $\Pi^{(1)}$}
% ==============================================================================

In Chapter~\ref{ch:methods}, we introduced the rank-truncated projected state $\rho^{(m)}_\theta = \Pi^{(m)}\rho_\theta\Pi^{(m)}$ as an abstract, mathematical construction to isolate the dominant eigenspace of a high-dimensional mixed state and construct dynamics-agnostic metrological bounds. We now demonstrate that spatial-mode demultiplexing (SPADE) acts as a physical, analog hardware-level implementation of this exact truncation scheme on a two-dimensional spatial subspace.

Recall the single-photon spatio-temporal mixed state $\rho_d$ in the image plane of the telescope, defined in Eq.~\eqref{eq:optical_density_matrix_chap}. In astronomical observations under extreme high-contrast regimes ($\epsilon \ll 1$), the overwhelming starlight acts as a massive coherent background that dwarfs the exoplanet's signal. The rank-truncation operator with $m=1$ provides the exact physical filter required to isolate and discard this uninformative starlight background.

Using the first two Hermite-Gaussian (HG) spatial modes defined in Eqs.~\eqref{eq:hg00_def}--\eqref{eq:hg10_def}, we recall the orthogonal projection operators:
\begin{align}
    \Pi_0 & = |HG_{00}\rangle\langle HG_{00}| \quad \text{(The Dominant Star Subspace)},                                         \\
    \Pi_1 & = |HG_{01}\rangle\langle HG_{01}| + |HG_{10}\rangle\langle HG_{10}| \quad \text{(The Perturbative Planet Subspace)}.
\end{align}
Under this mapping, the physical projection of the optical field onto the asymmetric modes is mathematically equivalent to applying the rank-truncation projector:
\begin{equation}
    \rho_d^{(1)} = \Pi^{(1)}\rho_d\Pi^{(1)}, \quad \text{where} \quad \Pi^{(1)} = \Pi_0 + \Pi_1.
    \label{eq:physical_rank_truncation}
\end{equation}
By routing the starlight into $\Pi_0$, the parameter estimation of the transverse separation $d$ is evaluated exclusively in the orthogonal subspace $\Pi_1$, shielding the weak signal from the starlight's shot noise. This is the exact optical analogue of our variational state-diagonalization protocol where the dominant eigenvalue is mapped to a known basis state to inspect the remaining subnormalized spectrum.

\subsection{Telescope Alignment as VQSE Majorization}
Before the physical rank-truncation projector can be executed, the optical system must be physically aligned with the bright star. In Section~\ref{sec:astrophysical_metrology_spade}, we detailed the star-centric alignment protocol, where the telescope's spatial coordinates are tuned to maximize the single-photon count $p_0$ in the fundamental mode $|HG_{00}\rangle$ (Eq.~\eqref{eq:p0_alignment}).

We now establish a formal bijection between this physical alignment and the mathematical majorization mechanism of VQSE. Recall from Chapter~\ref{ch:methods} that VQSE diagonalizes an unknown state by minimizing the energy expectation value of a sorted, non-degenerate local Hamiltonian $H_L = \sum_k E_k |e_k\rangle\langle e_k|$ (Eq.~\eqref{eq:vqse_cost_opt}). Because the true eigenvalues $\boldsymbol{\lambda}$ majorize the diagonal probability elements $\mathbf{p}$ in any basis ($\boldsymbol{\lambda} \succ \mathbf{p}$), minimizing the cost function $C(\beta)$ mathematically rotates the dominant eigenvector of $\rho$ into the lowest energy basis state $|e_1\rangle$.

In the SPADE pipeline, the physical star-centric alignment acts as a physical, analog realization of the VQSE majorization loop:
\begin{enumerate}
    \item **The State Mapping:** The mixed state $\rho_d$ acts as the unknown input state to be diagonalized.
    \item **The Hamiltonian Mapping:** We define a physical "Hamiltonian" in the HG basis:
          \begin{equation}
              H_{\text{align}} = 0 \cdot |HG_{00}\rangle\langle HG_{00}| + 1 \cdot |HG_{01}\rangle\langle HG_{01}| + 1 \cdot |HG_{10}\rangle\langle HG_{10}|.
          \end{equation}
    \item **The Cost Function Equivalence:** Minimizing the expectation value of $H_{\text{align}}$ corresponds to:
          \begin{equation}
              C_{\text{align}} = \operatorname{Tr}[H_{\text{align}}\rho_d] = p_1 = 1 - p_0.
          \end{equation}
\end{enumerate}
Maximizing the fundamental mode occupancy $p_0$ is mathematically identical to minimizing the energy cost $C_{\text{align}}$. The physical alignment knobs (fine-tuning the translation stages) serve as the physical "parameters" $\beta$, and the optical apparatus itself acts as the analog processor. By maximizing $p_0$, the physical alignment rotates the dominant "starlight" eigenvector of the optical field directly into the $|HG_{00}\rangle$ basis state.

% ==============================================================================
\section{Rigorous Proof of Variational TQFI Saturation}
% ==============================================================================

This physical correspondence can be formalized mathematically by proving that the classical Fisher Information extracted by SPADE perfectly saturates the ultimate variational Truncated QFI (TQFI) lower bound developed in our NISQ metrology track.

Recall the general definition of the TQFI for a subnormalized projected state under the generalized subnormalized fidelity $F_*$, established in Eq.~\eqref{eq:tqfi_bounds}:
\begin{equation}
    \mathcal{I}_*\left(d; \rho^{(m)}_d\right) \coloneqq 8 \lim_{\delta \to 0} \frac{1 - F_*\left(\rho^{(m)}_d, \rho^{(m)}_{d+\delta}\right)}{\delta^2}.
    \label{eq:tqfi_general_def}
\end{equation}
We evaluate this bound for a truncation rank of $m=1$ on the single-photon density matrix $\rho_d$ of the optical field (Eq.~\eqref{eq:optical_density_matrix_chap}), isolating the dominant star component. In this subnormalized projected representation, the physical distinguishability between $\rho_d$ and the shifted state $\rho_{d+\delta}$ is governed entirely by the transfer of probability amplitude from the fundamental mode $\Pi_0$ to the asymmetric mode $\Pi_1$.

The generalized subnormalized fidelity $F_*$ between these subnormalized operators is defined as:
\begin{equation}
    F_*(\sigma, \tau) \coloneqq \operatorname{Tr}\left[\sqrt{\sqrt{\sigma}\tau\sqrt{\sigma}}\right] + \sqrt{(1 - \operatorname{Tr}[\sigma])(1 - \operatorname{Tr}[\tau])}.
    \label{eq:f_star_definition}
\end{equation}
For $m=1$, the subnormalized projected operators are given by:
\begin{align}
    \rho_d^{(1)}          & = \lambda_0 |HG_{00}\rangle\langle HG_{00}|,    \\
    \rho_{d+\delta}^{(1)} & = \lambda_0' |HG_{00}'\rangle\langle HG_{00}'|,
\end{align}
where $\lambda_0 = p_0(d)$ and $\lambda_0' = p_0(d+\delta)$ are the fundamental mode occupancy probabilities.

Evaluating the trace term in Eq.~\eqref{eq:f_star_definition}:
\begin{equation}
    \operatorname{Tr}\left[\sqrt{\sqrt{\rho_d^{(1)}}\rho_{d+\delta}^{(1)}\sqrt{\rho_d^{(1)}}}\right] = \sqrt{p_0(d) p_0(d+\delta)} \left|\langle HG_{00} | HG_{00}' \rangle\right|.
\end{equation}
Using the wavefunctions and the overlap of the shifted spatial states, the inner product between the fundamental modes under an infinitesimal shift $\delta$ is:
\begin{equation}
    \left|\langle HG_{00} | HG_{00}' \rangle\right| \approx 1 - \frac{\eta\epsilon}{8w_0^2}\delta^2.
\end{equation}
Substituting the probability expressions and expanding the square of the generalized subnormalized fidelity to second order in $\delta$ yields:
\begin{equation}
    F_*\left(\rho^{(1)}_d, \rho^{(1)}_{d+\delta}\right)^2 \approx 1 - \frac{\eta \epsilon}{4w_0^2}\delta^2.
    \label{eq:fidelity_taylor_expansion_link}
\end{equation}
Substituting this Taylor expansion back into the TQFI definition (Eq.~\eqref{eq:tqfi_general_def}) yields:
\begin{equation}
    \mathcal{I}_*\left(d; \rho^{(1)}_d\right) = 8 \left( \frac{\eta \epsilon}{8w_0^2} \right) = \frac{\eta \epsilon}{w_0^2}.
    \label{eq:tqfi_result_saturated}
\end{equation}

This is a profound physical result. Comparing Eq.~\eqref{eq:tqfi_result_saturated} with the classical Fisher Information extracted by counting photons in the Hermite-Gaussian modes (Eq.~\eqref{eq:cfi_evaluation} / Eq.~\eqref{eq:spade_fisher_constant}), we find:
\begin{equation}
    \mathcal{I}_{*\text{, TQFI}}\left(d; \rho^{(1)}_d\right) \equiv \mathcal{I}_{\text{class, SPADE}}(d) = \frac{\eta \epsilon}{w_0^2}.
    \label{eq:saturation_equivalence}
\end{equation}
This equivalence proves that **the TQFI bounds we optimized variationally on NISQ simulators are the exact theoretical limits governing SPADE superresolution**. The classical measurement outcomes of SPADE completely saturate the ultimate quantum precision limit of the subnormalized projected space.

% ==============================================================================
\section{Dimensionality Sufficiency via Krylov Complexity Closure}
% ==============================================================================

A fundamental question in quantum metrology is whether a low-dimensional projective measurement is sufficient to capture the full metrological information of a high-dimensional system, or if full quantum state tomography is required. We resolve this question by linking the efficiency of SPADE to the **Krylov subspace convergence of Quantum Fisher Information** developed in Chapter~\ref{ch:methods} \cite{KrylovQFI2026}.

Recall that the exact QFI can be evaluated iteratively by generating the Krylov subspace of the state-dependent Liouvillian superoperator $\mathcal{K}_\rho$ acting on the parameter derivative operator $\mathcal{O}_0 = \partial_d \rho_d$ \cite{KrylovQFI2026}:
\begin{equation}
    \mathcal{K}_k\left(\mathcal{K}_\rho, \mathcal{O}_0\right) = \operatorname{span}\left\{\mathcal{O}_0, \mathcal{K}_\rho(\mathcal{O}_0), \dots, \mathcal{K}_\rho^{k-1}(\mathcal{O}_0)\right\}.
    \label{eq:krylov_subspace_link}
\end{equation}
Because the single-photon mixed state $\rho_d$ is highly rank-deficient and operates in the weak-field, faint-signal limit ($\eta \ll 1$) \cite{Santamaria2024single}, its state-dependent operator geometry is extremely constrained.

The Krylov distribution demonstrates that the Krylov complexity associated with the spatial displacement generator acting on $\rho_d$ closes at an exceptionally low dimension:
\begin{equation}
    k^* \le 2.
    \label{eq:krylov_closure}
\end{equation}
This mathematical closure proves that **full quantum state tomography is mathematically redundant**. Because the operator growth of the spatial derivative is confined to a 2-dimensional Krylov subspace, a simple 2-basis $HG_{00}/HG_{01}$ demultiplexer is guaranteed to capture nearly $100\%$ of the exact QFI \cite{KrylovQFI2026}. This provides the ultimate theoretical justification for the metrological efficiency of low-rank structured physical receivers.

% ==============================================================================
\section{Denoising as Polar Completion of Projected Operators}
% ==============================================================================

In Section~\ref{sec:astrophysical_metrology_spade}, we introduced the Physics-Informed Convolutional Neural Network (CNN) as a variational decoder to recover the separation $d$ under severe background noise, photon loss, and centroid drift. We now show that this machine learning recovery scheme is the exact physical analogue of the **polar completion** used to repair truncated VQSE updates on NISQ simulators.

On a quantum simulator, applying an unconstrained projection operator $P_M$ to a physical density matrix causes state leakage and non-unitarity. To maintain physical consistency, the polar completion constructs the unique nearest isometry $\widetilde{K}$ to the projected update in the Frobenius norm:
\begin{equation}
    \widetilde{K}^{(x)}_{y,\eta} = \bar{K}^{(x)}_{y,\eta} G_x^{-1/2}, \quad \text{where} \quad G_x = \sum_{y,\eta} \bar{K}^{(x)\dagger}_{y,\eta} \bar{K}^{(x)}_{y,\eta}.
    \label{eq:polar_completion_link}
\end{equation}

The Physics-Informed CNN acts as the physical, numerical analogue of this mathematical completion:
\begin{enumerate}
    \item **The Non-Unitary Projection:** Physical aberrations, spatial crosstalk ($\chi = 0.0035$), and photon losses act as a non-unitary projection that corrupts and leaks the spatial mode intensities.
    \item **The Isometric Reconstruction:** The CNN is trained with the physics-informed ratio of even HG modes $r_{20} \coloneqq p_2/p_0$ (which is structurally immune to symmetric alignment drift). By mapping the corrupted, non-unitary intensity profiles back to the clean, parameter-decoupled spatial separation $d$, the network numerically learns the unique isometric transformation that preserves information conservation.
\end{enumerate}
This allows the variational decoder to reconstruct the sub-wavelength separation with $>82\%$ fidelity, demonstrating how machine learning decoders physically execute the mathematical repair of projected quantum channels.

% ==============================================================================
\section{Global Precision Limits: QFIM Coordinate Transformation and Decoupling}
% ==============================================================================

To benchmark our 2-basis projection against the global quantum limits of multiparameter estimation, we analyze the joint Quantum Fisher Information Matrix (QFIM) $Q(s_0, d, \epsilon)$ derived by \v{R}eha\v{c}ek et al. \cite{Rehacek2017multiparameter}:
\begin{equation}
    Q(s_0, d, \epsilon) = 4 \begin{pmatrix}
        p_2 + 4\epsilon(1-\epsilon)P^2 & (\epsilon - 1/2)p_2 & -iwP                                   \\
        (\epsilon - 1/2)p_2            & p_2/4               & 0                                      \\
        -iwP                           & 0                   & 1 - \frac{w^2}{4 \epsilon(1-\epsilon)}
    \end{pmatrix},
    \label{eq:qfim_rehacek_link}
\end{equation}
where $s_0$ represents the midpoint between the sources, and $p_2, P, w$ are positive parameters associated with the Gaussian PSF profile.

\subsection{Coordinate Transformation to Star-Centric Frame}
In our experiment, the telescope is aligned to the brighter star (Source A) at position $x = s_0 - d/2$. To evaluate how information behaves under this star-centric alignment, we apply a linear coordinate transformation matrix $R$:
\begin{equation}
    \begin{pmatrix} x \\ d \\ \epsilon \end{pmatrix} = R \begin{pmatrix} s_0 \\ d \\ \epsilon \end{pmatrix}, \quad \text{where} \quad R = \begin{pmatrix} 1 & -1/2 & 0 \\ 0 & 1 & 0 \\ 0 & 0 & 1 \end{pmatrix}.
    \label{eq:coord_map_link}
\end{equation}

The covariance matrix of the star-centric frame $\Sigma(x, d, \epsilon)$ is mapped via $\Sigma(x, d, \epsilon) = R \Sigma(s_0, d, \epsilon) R^T$. Inverting this relationship, the transformed QFIM is given by:
\begin{equation}
    Q(x, d, \epsilon) = (R^T)^{-1} Q(s_0, d, \epsilon) R^{-1} = 4 \begin{pmatrix}
        p_2 + 4\epsilon(1-\epsilon)P^2   & (p_2 + 2(1-\epsilon)P^2)\epsilon  & -iwP                                  \\
        (p_2 + 2(1-\epsilon)P^2)\epsilon & (p_2 + (1-\epsilon)P^2)\epsilon^2 & -iwP\epsilon/2                        \\
        -iwP                             & -iwP\epsilon/2                    & 1 - \frac{w^2}{4\epsilon(1-\epsilon)}
    \end{pmatrix}.
    \label{eq:transformed_qfim_link}
\end{equation}

\subsection{Separation of Scale and Parameter Decoupling}
The matrix element $Q_{11}$ (associated with the star position $x$) is of order $\mathcal{O}(1)$, while $Q_{22}$ (separation $d$) is of order $\mathcal{O}(\epsilon)$. This disparity establishes a separation of scale between the alignment and the separation parameters.

Assuming the intensity ratio $\epsilon$ is known, we evaluate the sub-matrix $Q(x,d)$:
\begin{equation}
    Q(x, d) = 4 \begin{pmatrix}
        p_2 + 4\epsilon(1-\epsilon)P^2   & (p_2 + 2(1-\epsilon)P^2)\epsilon \\
        (p_2 + 2(1-\epsilon)P^2)\epsilon & \epsilon(p_2 + (1-\epsilon)P^2)
    \end{pmatrix}.
\end{equation}

In the presence of non-zero off-diagonal correlations, the optimal Cram\'{e}r-Rao variance is obtained by inverting the full QFI matrix:
\begin{equation}
    \Sigma^2_d = [Q(x,d)^{-1}]_{2,2} = \frac{1}{4\epsilon} \frac{1}{p_2 + (1-\epsilon)P^2} \left( 1 + \frac{\epsilon}{1-\epsilon} \frac{(p_2 + 2(1-\epsilon)P^2)^2}{p_2(p_2 + P^2)} \right).
    \label{eq:crb_full_link}
\end{equation}

As the relative intensity ratio $\epsilon \to 0$, the correlation term proportional to $\frac{\epsilon}{1-\epsilon}$ in Eq.~\eqref{eq:crb_full_link} vanishes. Thus, the alignment error and the separation estimation become completely decoupled:
\begin{equation}
    \lim_{\epsilon \to 0} \Sigma^2_d = \tilde{\Sigma}^2_d = \frac{1}{4\epsilon (p_2 + (1-\epsilon)P^2)}.
    \label{eq:decoupled_crb_link}
\end{equation}
This mathematically justifies our treating of the alignment procedure and the parameter estimation as independent subroutines, providing a rigorous information-theoretic foundation for our experimental tabletop analysis.

% ==============================================================================
\section{Experimental Calibration and Regression Synthesis}
% ==============================================================================

To physically verify the decoupled sub-Rayleigh precision limits derived above, we synthesize the results from the tabletop experimental configuration described in Section~\ref{sec:astrophysical_metrology_spade} using the physical noise parameterizations.

\subsection{Regression Calibration Models}
Recall that the average detected fraction $f_1 = n_1/(\eta n)$ incorporates physical cross talk ($\chi = 0.0035$) and stochastic misalignment into the effective noise parameter $\sigma_{\text{eff}}^2 \coloneqq \sigma_M^2 + 4w_0^2\chi$ (Eq.~\eqref{eq:sigma_effective}).

To construct our experimental calibration curves, data were collected by translating the B beam between $-200\,\mu\text{m}$ and $200\,\mu\text{m}$ ($d_a = d/w_0 \in [-0.67, 0.67]$). For each step, we recorded the combined photon counts $n_1 = C^n_{HG01} + C^n_{HG10}$ during $1\text{ s}$ of acquisition time, composed of $N_m = 100$ independent measures of $10\text{ ms}$ each.

\begin{figure}[h!]
    \centering
    \begin{center}
        \textbf{Calibration Curves for Separation $d_a$ at Four Intensity Ratios $\epsilon$:} \\
        \begin{tabular}{cc}
            $\epsilon = 0.018: \quad n_1^F(d_a) = a_1 + b_1 d_a^2$ & $\epsilon = 0.042: \quad n_1^F(d_a) = a_1 + b_1 d_a^2$ \\
            $\epsilon = 0.069: \quad n_1^F(d_a) = a_1 + b_1 d_a^2$ & $\epsilon = 0.096: \quad n_1^F(d_a) = a_1 + b_1 d_a^2$
        \end{tabular}
    \end{center}
    \caption{Quadratic regression curves of the first-order mode counts $n_1(d_a)$ versus dimensionless separation $d_a$ for the four experimental intensity ratios $\epsilon \in \{0.018, 0.042, 0.069, 0.096\}$. The fitting parameters $a_1$ and $b_1$ account for the physical cross talk ($\chi = 0.0035$) and dark count rates.}
    \label{fig:calibration_da_link}
\end{figure}

The fitting parameter $a_1$ corresponds to the baseline count rate from the effective noise parameter $\sigma_{\text{eff}}^2$, while $b_1$ represents the parameter-dependent signal scaling:
\begin{align}
    a_1 & \propto \eta \sigma_{\text{eff}}^2 = \eta (\sigma_M^2 + 4w_0^2\chi), \\
    b_1 & \propto \eta \epsilon w_0^2.
\end{align}
The resulting experimental uncertainties $\Delta d_a$ obtained via these quadratic calibrations show that our SPADE tabletop setup successfully beats the direct imaging theoretical limits for separations down to $d_a \approx 0.15$.

For the relative intensity ratio $\epsilon$, linear calibrations were performed at the four fixed separations $d_a \in \{0.2, 0.27, 0.47, 0.67\}$ established in Eq.~\eqref{eq:epsilon_cal_points}:

\begin{figure}[h!]
    \centering
    \begin{center}
        \textbf{Calibration Curves for Intensity $\epsilon$ at Four Separations $d_a$:} \\
        \begin{tabular}{cc}
            $d_a = 0.2: \quad n_1^F(\epsilon) = \alpha_1 + \beta_1 \epsilon$  & $d_a = 0.27: \quad n_1^F(\epsilon) = \alpha_1 + \beta_1 \epsilon$ \\
            $d_a = 0.47: \quad n_1^F(\epsilon) = \alpha_1 + \beta_1 \epsilon$ & $d_a = 0.67: \quad n_1^F(\epsilon) = \alpha_1 + \beta_1 \epsilon$
        \end{tabular}
    \end{center}
    \caption{Linear regression calibration curves of the first-order mode counts $n_1(\epsilon)$ versus relative intensity $\epsilon$ for the four fixed dimensionless separations $d_a \in \{0.2, 0.27, 0.47, 0.67\}$.}
    \label{fig:calibration_eps_link}
\end{figure}

The fitting parameter $\alpha_1$ accounts for the background counts from spatial crosstalk, while $\beta_1$ scales linearly with the separation $d_a^2$. The resulting experimental uncertainties $\Delta\epsilon$ match our theoretical stochastic error models (Eq.~\eqref{eq:propagation_uncertainty_d}--\eqref{eq:std_dev_fraction}) within a factor of two, verifying the metrological advantage of SPADE in realistic high-contrast regimes.